% Fuente: definiciones de momentos y usos desarrollados en EQT y Discusión.
\begin{frame}[label=nav-frame-099]{Momentos de $Q_\beta$: cálculo y uso posterior}
\TablaSetup
\begin{tabular}{@{}p{.19\textwidth}p{.30\textwidth}p{.44\textwidth}@{}}
\TablaCabecera{Objeto} & \TablaCabecera{Derivada} & \TablaCabecera{Uso en la deducción} \\
\TablaLinea
$P_\beta^{abcd}$ & $\dfrac{\partial Q_\beta}{\partial R_{abcd}}$ & Contracción $\mathcal R^{(\beta)}_{ab}$, doble divergencia, frontera y contraste LL. \\[3mm]
$P^{(\beta)}_{ab}$ & $\dfrac{\partial Q_\beta}{\partial g^{ab}}$ & Ecuación métrica e identidad algebraica para $\mathcal R^{(\beta)}_{ab}$. \\[3mm]
$\Pi_\beta^a$ & $\dfrac{\partial Q_\beta}{\partial u_a}$ & Ecuación escalar, frontera y chequeo de Bianchi. \\[3mm]
$F_\phi^{(\beta)}$ & $\dfrac{\partial Q_\beta}{\partial\phi}$ & Dependencia explícita en $\phi$ y ecuación escalar. \\
\TablaLinea
\end{tabular}\par
\vspace{3mm}
\footnotesize
$u_a=\nabla_a\phi$, $X=u^a u_a$. Las derivadas mantienen fijos los
demás argumentos. Primero calculamos los momentos y luego los combinamos.
\end{frame}
